% Descent tools used before the effective-descent proof in Chapter 1.
\section{The geometric tool that makes object descent possible}
\label{desc:tools}

We now need a theorem that produces a scheme from local schemes.
Before applying it to our curves, we explain its input: a line bundle
that supplies compatible projective coordinates. Merely knowing that
every local curve admits \emph{some} projective embedding does not supply
compatible embeddings on overlaps. The particular line bundle below is
determined by the pointed curve itself, so every transition isomorphism
automatically acts on it.

\begin{Def}[The line bundle supplied by a stable pointed family]
\label{desc:canonical-bundle}
Start with the family $(p:C\to S,\sigma_1,\ldots,\sigma_n)$ of
Definition~\ref{def:stable-family}. Each marking has an image
$D_a=\sigma_a(S)$ in $C$. We form
\[
 D=D_1+\cdots+D_n,
 \qquad L_C=\omega_{C/S}(D)
       =\omega_{C/S}\otimes\mathcal O_C(D).
\]
Here $\omega_{C/S}$ is the \emph{relative dualizing sheaf}: the line
bundle playing the role of relative differential forms for a family
whose fibers are allowed to have nodes. The notation $(D)$ means
tensoring with the line bundle associated to the indicated divisor.
\end{Def}

This definition uses two results about families of schemes, which we
state as external inputs. A section through the smooth locus of a
separated family of curves is a relative effective Cartier divisor.
Concretely, its image is locally cut out by one equation that is a
non-zero-divisor, and this description remains a Cartier divisor after
base change. One can check this in a relative smooth coordinate, where
the section has equation $z=0$; the general statement follows by
\'etale localization. See Stacks Project, Section 44.3,
Lemma 44.3.1, and Section 31.23, Lemma 31.23.8
\cite{Stacks}. Thus all the $\mathcal O_C(D_a)$ above are defined and
commute with base change.

The second input is the relative duality theorem in its curve form:
if $p:C\to S$ is proper, flat, and finitely presented, with
Gorenstein geometric fibers of pure dimension one, then
$\omega_{C/S}$ is invertible and, for every $b:T\to S$, there is a
canonical isomorphism
\begin{equation}\label{desc:dualizing-basechange}
 \operatorname{pr}_C^*\omega_{C/S}
       \simeq \omega_{C\times_S T/T}.
\end{equation}
These isomorphisms respect further base change and isomorphisms of
families. This is a theorem about proper flat families; it does not
require us to know that our moduli category is algebraic. A precise
curve formulation is Stacks Project, \MCtag{0E6R},
Lemma 109.19.3, with its duality references in Section 109.19
\cite{Stacks}. We use this duality statement as a black box.

Why does it apply? A smooth curve has regular local rings; a node is
locally a hypersurface singularity, with completed local ring
$\overline{k}[[x,y]]/(xy)$. These rings are Gorenstein: their dualizing
modules have rank one and are free. Consequently the geometric fibers
in our definition satisfy the theorem. On the smooth locus,
$\omega_{C/S}$ agrees with relative differentials. On a single nodal
fiber its sections can be described on the normalization as
differentials with at most simple poles over each node, with the two
residues summing to zero. This explains why ordinary regular
differentials on a smooth curve need a replacement at a node.

\begin{Lem}[Functorial relative ampleness]
\label{desc:ample-canonical}
For a stable pointed family, $L_C=\omega_{C/S}(D)$ is relatively ample
over $S$. It commutes canonically with arbitrary base change. A pointed
isomorphism $a:C\xrightarrow{\sim}C'$ over $S$ therefore gives a
canonical isomorphism $a^*L_{C'}\simeq L_C$, compatible with
composition of pointed isomorphisms.
\end{Lem}

\noindent\emph{Verification.}
The base-change assertion follows from
\eqref{desc:dualizing-basechange} and from the relative Cartier divisor
property of the sections. For ampleness, restrict to a geometric fiber
and normalize any one of its irreducible components. If that
normalization has genus $h$ and $r$ distinguished points, counting the
branches of nodes and the markings, then
\begin{equation}\label{eq:component-degree}
 \deg L_C|_{\widetilde E}=2h-2+r>0
\end{equation}
by stability. Positivity on every component implies
ampleness on the fiber. The precise curve criterion is Stacks Project,
\MCtag{0B5Y}, Lemma 33.44.15; degrees are unchanged on normalization
by Lemma 33.44.4. For a proper family, ampleness on a fiber implies
relative ampleness on a neighborhood of its base point
(\MCtag{0D2S}, Lemma 37.50.3). These neighborhoods cover $S$, which
proves relative ampleness. The final assertion follows because an
isomorphism carries each labelled section to the corresponding labelled
section, and the duality identifications are functorial.\hfill$\square$

In ordinary language, every stable pointed family comes with the
\emph{same rule} for choosing an ample line bundle. We have not added a
line bundle to the objects of $\Mgn$: we have constructed one from data
already present. For a smooth unpointed genus-two curve $C$, this rule
is simply $L_C=\omega_C$, of degree $2$. For a genus-two curve made from
two smooth elliptic curves joined at one node, its degree on each
component is $1$. The positivity check works on both components even
though the curve is singular.

\begin{Th}[Effective descent with a compatible ample line bundle]
\label{desc:polarized-theorem}
Let $\{S_i\to S\}$ be an \textup{\'etale} cover of a scheme.
Suppose that for every $i$ we are given a proper, flat, finitely
presented morphism $p_i:Y_i\to S_i$ and a relatively ample invertible
sheaf $L_i$ on $Y_i$. Suppose also that there are isomorphisms
\[
 a_{ij}:Y_i|_{S_{ij}}\xrightarrow{\sim}Y_j|_{S_{ij}},
 \qquad
 \lambda_{ij}:a_{ij}^*(L_j|_{S_{ij}})
                   \xrightarrow{\sim}L_i|_{S_{ij}},
\]
where $S_{ij}=S_i\times_S S_j$. Assume the scheme isomorphisms satisfy
$a_{ik}=a_{jk}\circ a_{ij}$ on $S_{ijk}$, and the line-bundle
isomorphisms satisfy
\begin{equation}\label{desc:bundle-cocycle}
 \lambda_{ik}=\lambda_{ij}\circ a_{ij}^*\lambda_{jk}
 \quad\text{on }S_{ijk},
\end{equation}
using the canonical identifications of iterated pullbacks.
Then there are a proper, flat, finitely presented scheme $Y\to S$, a
relatively ample invertible sheaf $L$ on $Y$, and isomorphisms
$(Y,L)|_{S_i}\simeq(Y_i,L_i)$ recovering the prescribed descent data.
The same statement holds for fpqc covers.
\end{Th}

\noindent\textbf{Source and proof strategy.}
The scheme-and-line-bundle effectivity theorem is SGA~1,
Expos\'e VIII, Proposition 7.8, p.~174 of the recomposed edition
\cite{SGA1}. Its argument descends a graded algebra and uses
relative $\Proj$. Vistoli's later exposition \cite{Vistoli},
Section 4.3.3, Theorem 4.38 and proof, pp.~96--103, gives the detailed
version for a functorial ample bundle on a class of proper flat
finitely presented families. His theorem applies this method to a
whole moduli problem; the statement above isolates the descent datum
that we need. Neither source says that arbitrary projective schemes
descend effectively without compatible polarization data.

\noindent\textbf{Our guided explanation of the construction.}
The input is local spaces together with instructions for identifying
them and their line bundles on overlaps. The output must be one
scheme over $S$, with identifications that reproduce precisely those
instructions. We describe the \emph{\'etale} case used in these notes.

First work over an affine open of $S$. Refine the cover by affine open
subschemes of its members. Because \'etale maps are open and $S$ is
quasi-compact, finitely many of their images cover $S$. Their disjoint
union $U$ is affine, and $U\to S$ is faithfully flat and
quasi-compact. The map need not be finite: it is the \emph{list of
affine covering schemes} that is finite. All the original data restrict
to $U$ and $U\times_S U$. After carrying out descent there, the
uniqueness of morphism descent identifies the answer with the data on
every original $S_i$. Answers over different affine opens of $S$ glue
by ordinary Zariski gluing. This reduction justifies using one uniform
power of a line bundle in the following steps.

\begin{enumerate}
\item \emph{Turn the local families into closed subschemes of projective
bundles.}
Write $p_U:Y_U\to U$ and $L_U$ for the disjoint union of the refined
data. Relative Serre vanishing and the relative very-ampleness theorem
allow us to choose $m>0$ such that $L_U^{\otimes m}$ is relatively very
ample and $H^j(Y_u,L_u^{\otimes m})=0$ for every $j>0$ on every fiber. Cohomology and
base change then makes
\[
 E_U=(p_U)_*L_U^{\otimes m}
\]
a finite locally free sheaf, with formation commuting with base change.
These are substantial scheme-theoretic inputs, not consequences of the
stack axioms. The cohomology statement, including arbitrary bases, is
Vistoli, Proposition 4.37, pp.~95--96; the passage to a sufficiently
large power occurs in the proof of Theorem 4.38. The evaluation map
$p_U^*E_U\to L_U^{\otimes m}$ gives a closed immersion
\[
 Y_U\hookrightarrow\PP_U(E_U)
           :=\Proj_U\bigl(\operatorname{Sym}^{\bullet}E_U\bigr).
\]
Here relative $\Proj$ is the construction that turns a graded algebra
over a base into a projective scheme over that base.

\item \emph{Descend the projective coordinates.}
The isomorphisms $a_{ij}$ and $\lambda_{ij}$ induce compatible
isomorphisms between the pullbacks of $E_U$ on $U\times_S U$.
Equation~\eqref{desc:bundle-cocycle} gives their cocycle condition.
Faithfully flat descent for quasi-coherent modules now constructs a
sheaf $E$ on $S$ whose pullback is $E_U$. Local freeness and finite rank
can be checked faithfully flat locally, so $E$ is a vector bundle.
We use Vistoli, Theorem 4.23, pp.~85--90, for this descent theorem.
It is the algebraic input that replaces elementary gluing of vector
bundles on open subsets.

Thus the projective bundles are all pullbacks of the single scheme
$\PP_S(E)$. Functoriality of evaluation ensures that the given
isomorphisms identify the embedded copies of $Y_U$ in this common
projective bundle. This is the step for which compatible line bundles
were essential.

\item \emph{Descend the equations.}
Let $\mathcal I_U\subset\mathcal O_{\PP_U(E_U)}$ be the ideal sheaf
of the embedded $Y_U$. The compatibility just checked gives this ideal
sheaf a descent datum over the fpqc cover
$\PP_U(E_U)\to\PP_S(E)$. Quasi-coherent descent constructs an ideal
sheaf $\mathcal I\subset\mathcal O_{\PP_S(E)}$; both its inclusion
and its closure under multiplication can be checked after the cover.
Let $Y\subset\PP_S(E)$ be the closed subscheme cut out by
$\mathcal I$. Formation of this closed subscheme commutes with flat
base change, so $Y\times_S U\simeq Y_U$ with the \emph{specified}
transition identifications. This is the projective-bundle form of the
graded-algebra/$\Proj$ construction in SGA~1.

\item \emph{Recover the line bundle and check the family.}
Now that $Y$ exists, the $L_i$ are line bundles on a cover of $Y$ with
compatible transition maps. Quasi-coherent descent constructs $L$;
invertibility is checked on that cover. Properness, flatness, and finite
presentation of $Y\to S$ follow from their fpqc locality on the base
(Stacks Project, \MCtag{02L1}, \MCtag{02L2}, \MCtag{02L0},
Lemmas 35.23.16, 35.23.17, 35.23.15). Relative ampleness of $L$
can be checked on geometric fibers and then near each base point
using the proper-family ampleness theorem above.
\end{enumerate}

This explains why the theorem returns a \emph{scheme}. It does not
begin by assuming that the sheaf quotient of the $Y_i$ is a scheme.
The vector bundle and the descended equations construct that scheme.
The theorem has not yet supplied markings or a map to a target; we
will descend those additional morphisms after applying it.

\begin{Ej}\label{desc:exercise-polarization}
Suppose the local pointed families have transition isomorphisms
$a_{ij}$ satisfying the cocycle condition. Explain why choosing
$L_i=\omega_{C_i/S_i}(\sum_a\sigma_{a,i})$ supplies the
$\lambda_{ij}$ and equation~\eqref{desc:bundle-cocycle} without an
additional choice. Why would choosing an unrelated ample line bundle
on each $C_i$ not by itself do this?
\end{Ej}
